\section{Methodology}\label{methodology}

This section describes the two primary identification strategies---firm
fixed-effects OLS and Double Machine Learning---as well as the
Fama-MacBeth procedure used for temporal and cross-sectional analysis,
the cap-group classification, and the long-short portfolio construction.

\subsection{Firm Fixed-Effects OLS}\label{firm-fixed-effects-ols}

The baseline identification strategy is a within-firm panel regression.
Let \(Y_{\text{it}}\) denote the open-to-close price change
(\texttt{return}) for firm \(i\) on day \(t\), \(D_{\text{it}}\) the
sentiment treatment, and \(X_{\text{it}}\) a vector of controls. The
estimating equation is:

\[Y_{\text{it}} = \alpha_{i} + \beta D_{\text{it}} + \mathbf{\gamma}^{'}X_{\text{it}} + \varepsilon_{\text{it}},\]

where \(\alpha_{i}\) is a firm fixed effect and
\(\varepsilon_{\text{it}}\) is an idiosyncratic error. The parameter of
interest is \(\beta\), the within-firm partial effect of sentiment on
same-day return after conditioning on observable time-varying
confounders.

Firm fixed effects absorb all time-invariant, firm-specific factors that
could jointly influence sentiment coverage and price levels---including
industry membership, market capitalization quintile, analyst following,
and the firm's typical trading clientele. The identifying variation is
purely within-firm and over time. This is appropriate because
Bloomberg's Twitter sentiment measure has markedly heterogeneous
coverage across firms
(Section~\ref{measurement-limitations}): the fixed
effect controls for the fact that firms with high average sentiment may
simply be large, liquid, or widely-followed, and thus differ from
low-coverage firms on dimensions that also affect returns.

The control vector \(X_{\text{it}}\) in the full specification (Model~2)
includes: intraday price range (\texttt{px\_high}, \texttt{px\_low}),
market capitalization (\texttt{mkt\_cap}), book equity
(\texttt{total\_equity}), leverage (\texttt{debt\_to\_equity}), trading
volume (\texttt{volume}), news sentiment (\texttt{news\_sent}), negative
tweet count (\texttt{twitter\_neg\_count}), 30-day RSI (\texttt{rsi\_30}),
and the 50-day moving average (\texttt{ma\_50}).
Standard errors are heteroskedasticity-robust throughout.

Four treatment constructions are estimated. Model~1 uses the continuous
daily average sentiment score
\(D_{\text{it}} = \text{twitter\_sent} \in \lbrack - 1, + 1\rbrack\).
Model~2 augments Model~1 with the full control vector. Model~3 replaces
the treatment with a negative-clipped measure,
\(D_{\text{it}} = \text{min}(\text{twitter\_sent}_{\text{it}},\, 0)\),
which equals zero on positive-sentiment days and the raw score on
negative-sentiment days---isolating whether the return effect is
asymmetric in the negative tail. Model~4, following Teti et al.\ (2019),
uses daily negative tweet counts (\texttt{twitter\_neg\_count}) as the
treatment. All four models include firm fixed effects.

\textbf{Limitation.} The FE-OLS specification assumes that the
confounder--outcome and confounder--treatment relationships are both
linear. Nonlinear interactions among controls---for example, between
momentum (RSI) and sentiment during high-volatility regimes---will
contaminate the \(\beta\) estimate if the true data-generating process
is non-additive. The DoubleML specification in
Section~\ref{double-machine-learning} relaxes this
assumption.

\subsection{Double Machine Learning}\label{double-machine-learning}

The primary causal estimator is the Partially Linear Regression (PLR)
model of Chernozhukov et al.\ (2018), estimated via the \texttt{DoubleML}
Python package. The PLR model is:

\[\begin{matrix}
Y_{\text{it}} & = \theta_{0}D_{\text{it}} + g_{0}(X_{\text{it}}) + U_{\text{it}},\quad E\lbrack U_{\text{it}} \mid D_{\text{it}},X_{\text{it}}\rbrack = 0, \\
D_{\text{it}} & = m_{0}(X_{\text{it}}) + V_{\text{it}},\quad E\lbrack V_{\text{it}} \mid X_{\text{it}}\rbrack = 0, \\
\end{matrix}\]

where \(\theta_{0}\) is the causal parameter of interest, and
\(g_{0}( \cdot )\) and \(m_{0}( \cdot )\) are unknown nuisance functions
capturing the conditional expectations of the outcome and treatment
given the full confounder vector \(X_{\text{it}}\), respectively.
Critically, \(g_{0}\) and \(m_{0}\) are estimated
nonparametrically---here using XGBoost as the nuisance learner---so no
linearity or additive separability is assumed for the confounders.

\textbf{Cross-fitting.} To prevent overfitting bias from contaminating
the treatment effect estimate, I use \(K = 5\) folds of cross-fitting:
the nuisance functions \(\widehat{g}\) and \(\widehat{m}\) are trained
on four folds and evaluated on the held-out fold, so the residuals
\({\widehat{U}}_{\text{it}}\) and \({\widehat{V}}_{\text{it}}\) are
always out-of-sample predictions. The final coefficient
\(\widehat{\theta}\) is obtained by regressing
\({\widehat{U}}_{\text{it}}\) on \({\widehat{V}}_{\text{it}}\)---the
partialled-out outcome on the partialled-out treatment---aggregated
across folds. The primary treatment-effect estimates use a single
cross-fitting draw; standard errors are computed from the \texttt{DoubleML}
package's default asymptotic variance formula. For the lag-decay and
impact-persistence robustness specifications (\texttt{lag\_cates.py}),
I repeat the full cross-fitting procedure over 20 random fold assignments
and derive confidence intervals from 1,000 bootstrap draws, which
stabilizes coefficient estimates against fold-specific randomness in the
smaller effective samples obtained when conditioning on lagged or lead
return windows.

\textbf{XGBoost configuration.} The nuisance learner is XGBoost
(\texttt{XGBRegressor}), a gradient-boosted tree ensemble. Gradient
boosting is well-suited to this setting because it handles mixed-type
regressors, captures nonlinear and interaction effects implicitly
through tree splits, and is robust to the collinearity among sentiment,
momentum, and volume variables that can destabilize linear nuisance
models. The same confounder set used in the FE-OLS full model is
supplied to both nuisance learners.

\textbf{Treatment variants.} Three treatment series are estimated
separately in the DoubleML framework: (1) \texttt{twitter\_sent}
(continuous); (2) \texttt{twitter\_neg\_count} (count of
negative-polarity tweets, following Teti et al.\ (2019)); and (3)
\texttt{news\_sent} (continuous news sentiment). A fourth specification
estimates each treatment against \texttt{px\_high} (intraday price high)
as the outcome rather than \texttt{return}, to test whether sentiment's
effect on the intraday price ceiling differs from its effect on the
close-to-open spread.

\textbf{Identifying assumption.} Under the PLR, \(\theta_{0}\) is
identified if
\(E\lbrack U_{\text{it}} \mid D_{\text{it}},X_{\text{it}}\rbrack = 0\)---that
is, if the confounder vector \(X_{\text{it}}\) fully controls for the
non-random assignment of sentiment to firms and days. This is a
conditional unconfoundedness assumption. It does not require an
instrument. It does require that no omitted variable jointly drives both
\texttt{twitter\_sent} and \texttt{return} after conditioning on the
observable controls. The plausibility of this assumption is assessed
through the placebo test described in
Section~\ref{robustness-tests}.

\subsection{Fama-MacBeth Procedure and Temporal
Analysis}\label{fama-macbeth-procedure-and-temporal-analysis}

For the replication and temporal analysis, I follow Fama and MacBeth
(1973) and Gu and Kurov (2020) in estimating daily cross-sectional
regressions and aggregating the time series of coefficients. On each
trading day \(t\), I estimate:

\[r_{\text{it}} = a_{t} + b_{t} \cdot \text{twitter\_sent}_{i,t - 1} + \mathbf{c}_{t}^{'}X_{\text{it}} + e_{\text{it}},\]

where \(r_{\text{it}}\) is the Fama-French-Carhart four-factor-adjusted
return (\texttt{return\_oo\_adj}) for firm \(i\) on day \(t\), and the
treatment is the one-day lagged sentiment score. The confounder vector
\(X_{\text{it}}\) includes five lags each of \texttt{return\_oo\_adj},
\texttt{abnorm\_vol} (abnormal volume), \texttt{vol\_rs}
(Rogers-Satchell realized volatility), \texttt{log\_mkt\_cap}, and
\texttt{bid\_ask\_spread}. The Fama-MacBeth estimate of the sentiment
effect is \(\bar{b} = T^{- 1}\sum_{t = 1}^{T}{\widehat{b}}_{t}\), with
Newey-West standard errors computed over the time series of
\({\widehat{b}}_{t}\) using a bandwidth of five lags to correct for
serial correlation. Days on which fewer than 30 firms report non-missing
sentiment are excluded from the aggregation.

\textbf{Risk adjustment.} Factor-adjusted returns
(\texttt{return\_oo\_adj}) are constructed by estimating full-sample
Fama-French-Carhart four-factor betas for each firm and subtracting the
fitted four-factor component from the raw open-to-open return. Factor
data are obtained from Kenneth French's data library; the sample is
truncated at the last available factor date (February 28, 2026) to avoid
extrapolation. The residual \texttt{return\_oo\_adj} captures
idiosyncratic price variation that cannot be attributed to market, size,
value, or momentum exposures, and is the appropriate dependent variable
for tests of sentiment-specific information.

\textbf{Temporal windows.} The Fama-MacBeth procedure is applied
separately over annual, quarterly, and monthly windows to characterize
how the sentiment coefficient has evolved from 2016 to 2026. Results are
reported graphically as coefficient-and-confidence-band plots and, for
the annual cadence, in a tabular format. Monthly estimates are less
reliable due to small cross-sections in some months and are used for
visualization only.

\subsection{Cap-Group Classification and Russell 3000
Analysis}\label{cap-group-classification-and-russell-3000-analysis}

To test whether the sentiment effect is heterogeneous across the
market-capitalization spectrum, I classify all Russell 3000 firms into
three groups based on their time-averaged market capitalization:

\begin{itemize}
\item
  \textbf{Large cap:} \({\overline{\text{mkt\_cap}}}_{i} > \$ 10\)B
\item
  \textbf{Mid cap:} \(\$ 2\)B
  \(\leq {\overline{\text{mkt\_cap}}}_{i} \leq \$ 10\)B
\item
  \textbf{Small cap:} \({\overline{\text{mkt\_cap}}}_{i} < \$ 2\)B
\end{itemize}

I use time-averaged rather than daily market capitalization to assign
group membership, so that firms do not migrate between groups due to
price volatility within the sample window. This ensures that the
cap-group comparison is stable and that within-group temporal analysis
is not confounded by within-firm cap migration.

The Russell 3000 analysis proceeds in three stages. First, I estimate a
pooled Panel OLS with firm fixed effects that includes a
\texttt{twitter\_sent} \(\times\) \texttt{log\_mkt\_cap} interaction
term, testing whether the sentiment coefficient is linearly decreasing
in firm size across the full cross-section. Second, I run the
Fama-MacBeth procedure separately for each cap group and plot the annual
coefficient series. Third, I compute a daily long-short portfolio within
each cap group to assess the economic magnitude of the cross-sectional
dispersion in sentiment.

\subsection{Long-Short Portfolio
Construction}\label{long-short-portfolio-construction}

To assess the economic significance of the sentiment-return
relationship, I construct a daily long-short portfolio following Gu and
Kurov (2020). % [R2]: Clarify whether decile sort uses same-day (contemporaneous) or prior-day
% (lagged) sentiment. "Contemporaneous" is noted in the table caption but an
% open-to-open holding period implies prior-day sorting. Resolve before submission.
On each trading day \(t\), all firms with non-missing
\texttt{twitter\_sent} are ranked into deciles. The strategy takes a
long position in firms in the top decile (highest sentiment) and a short
position in firms in the bottom decile (lowest sentiment), holding for
24 hours on an open-to-open basis. Portfolio returns are the
equal-weighted average return across firms in each leg. For the
cap-group analysis, deciles are formed separately within each cap group,
so the strategy is self-contained per segment and not dominated by
large-cap firms with superior sentiment coverage.

Performance is summarized by annualized Sharpe ratio (assuming 252
trading days), annualized return, and win rate (fraction of days on
which the long leg outperforms the short leg). All reported figures are
before transaction costs, and the strategy should be interpreted as an
upper bound on achievable risk-adjusted performance.

\subsection{Robustness Tests}\label{robustness-tests}

Four robustness checks discipline the causal interpretation of the
primary estimates.

\textbf{Placebo test.} Yesterday's return (\texttt{lag1}) is regressed
on today's Twitter sentiment within firms. Under the null of no reverse
causality, the coefficient should be indistinguishable from zero. A
large, significant coefficient indicates that today's sentiment responds
to yesterday's price movement---feedback from prices to social media---and
constitutes evidence against a clean forward causal interpretation of the
primary estimates.

\textbf{No-reversal test.} I enter five sentiment lags jointly
(\(\text{twitter\_sent}_{i,t - 1}\) through
\(\text{twitter\_sent}_{i,t - 7}\), at lags 1, 2, 3, 5, 7) and test
whether only lag~1 is significant. Under a causal interpretation in
which information is absorbed within one trading day, the multi-lag
regression should show significance concentrated at lag~1 with
coefficients at longer lags near zero.

\textbf{Lag decay.} Each sentiment lag is estimated separately in a
univariate DoubleML specification with the same confounder set as the
primary model. If the coefficient is large at lag~1 and decays
monotonically to zero by lag~5 or lag~7, this is consistent with a
transient but real sentiment effect. A pattern of large, roughly equal
coefficients at all lags is inconsistent with causal transmission
through prices and more consistent with a common driver of both
sentiment and returns.

\textbf{Impact persistence.} Current-day sentiment is regressed on
future returns at leads \(+ 1\) through \(+ 7\) (same intervals as lag
decay). Under the hypothesis that positive sentiment today drives
positive returns tomorrow, the lead-1 coefficient should be positive and
significant, with the effect decaying at longer leads. A flat or
sign-inconsistent pattern across leads is evidence against a durable
predictive signal.
